\documentclass[11pt]{article}
\usepackage[margin=1in]{geometry}
\usepackage[T1]{fontenc}
\usepackage{amsmath,amssymb,booktabs,hyperref}
\usepackage[numbers]{natbib}
\hypersetup{colorlinks=true,citecolor=blue,urlcolor=blue}

\title{SAIS: Semantic Anchor \& Intent Score\\
\large A Fact-Aware, Structure-Aware Metric for Reference-Based Evaluation of Generated Text}
\author{Bharath Raghavan\\Independent Researcher}
\date{October 24, 2026}

\begin{document}
\maketitle

\begin{abstract}
Automatic metrics for generated text commonly emphasize lexical overlap or
semantic similarity, leaving factual completeness, unsupported factual
additions, structural ordering, and explicit requirements only indirectly
represented. This paper introduces the \emph{Semantic Anchor \& Intent Score}
(SAIS), a reference-based metric that combines a normalized embedding-based
meaning score, idea-block sequence alignment, a fact-level $F_\beta$ score,
an explicit constraint gate, and a symmetric length-ratio penalty. The design
makes factual and constraint failures multiplicative rather than compensable by
surface-level semantic similarity. We specify the metric, its implementation
interface, diagnostics, and a reproducible protocol for validation against
human judgments. SAIS is proposed as an evaluation framework; empirical
correlations are intentionally not claimed until they are measured on released,
human-scored benchmarks.
\end{abstract}

\section{Introduction}

The quality of a language-model response cannot be reduced to shared words.
Two valid responses may use different wording, while a fluent response may
silently omit a date, reverse a causal relation, introduce an unsupported
number, or violate a mandatory instruction. These failure modes are especially
important for summarization, question answering, and constrained generation.

BLEU~\citep{papineni2002bleu} measures modified $n$-gram precision with a
brevity penalty. ROUGE~\citep{lin2004rouge} measures overlap with a reference
and is widely used for summarization. Both are useful and inexpensive, but
their lexical basis can under-credit paraphrase and does not directly define
fact-level correctness. BERTScore~\citep{zhang2020bertscore} uses contextual
token embeddings and better captures semantic relatedness, but similarity alone
is not a sufficient test for required facts or explicit hard constraints.

SAIS is not proposed as a replacement for human evaluation. It makes the
evaluative assumptions explicit and returns component diagnostics, allowing
users to investigate a low score rather than treating it as an opaque label.

\section{Problem Definition}

Let $C$ be candidate text, $R$ a human-authored reference, and
$E_R,E_C$ the respective sets of salient facts. A fact may be a date, number,
named entity, identifier, or curated domain assertion. Let $\mathcal{K}$ be a
list of explicit constraints, such as a required statement, forbidden
disclosure, or formatting rule. SAIS returns a score in $[0,1]$ and a
diagnostic report containing matched facts, missing facts, candidate-only facts,
and constraint violations.

\section{The SAIS Formulation}

The overall score is
\begin{equation}
\operatorname{SAIS}(C,R)=
[\alpha S_{\mathrm{mesh}}(C,R)+(1-\alpha)S_{\mathrm{logic}}(C,R)]
S_{\mathrm{fact}}(C,R)\Gamma(C;\mathcal{K})BP(C,R),
\label{eq:sais}
\end{equation}
where $\alpha\in[0,1]$. Each factor is constrained to $[0,1]$. The
multiplicative terms reflect the intended policy: a candidate with a material
factual or hard-constraint failure should not receive a high final score merely
because its phrasing is semantically similar to the reference.

\subsection{Meaning Mesh}

Let $u=\operatorname{Embed}(C)$ and $v=\operatorname{Embed}(R)$ be normalized
sentence-level embeddings, for example from Sentence-BERT~\citep{reimers2019sentencebert}.
Because cosine similarity lies in $[-1,1]$, SAIS maps it to the unit interval:
\begin{equation}
S_{\mathrm{mesh}}(C,R)=\frac{1+\cos(u,v)}{2}.
\end{equation}
This component recognizes semantically equivalent paraphrases but is
deliberately insufficient on its own to establish factual faithfulness.

\subsection{Logic Flow}

The candidate and reference are segmented into ordered idea blocks
$B_C=(b_1,\ldots,b_m)$ and $B_R=(r_1,\ldots,r_n)$. For a block pair, let
$s(b_i,r_j)\in[0,1]$ be the normalized embedding cosine. Its position-aware
weight is
\begin{equation}
\omega(i,j)=\exp\left[-\lambda\left|\frac{i}{m}-\frac{j}{n}\right|\right],
\quad \lambda\geq0.
\end{equation}
The directional reference-coverage logic score is
\begin{equation}
S_{\mathrm{logic}}(C,R)=\frac{1}{n}\sum_{j=1}^{n}
\max_{i\in\{1,\ldots,m\}}\left[s(b_i,r_j)\omega(i,j)\right].
\label{eq:logic}
\end{equation}
Normalized positions are important: they compare relative location rather than
assuming the candidate and reference have identical numbers of blocks.
Production deployments may replace sentence segmentation with discourse or
proposition segmentation.

\subsection{Fact Anchor}

Facts require canonical comparison rather than raw string equality. The
canonicalization function normalizes case and whitespace, parses supported
dates (for example, ``October 24, 2026'' and ``2026-10-24''), normalizes
numbers, and may apply a domain-maintained alias table. Matches are one-to-one
to prevent duplicate credit.

Let $M=|E_R\cap_{\mathrm{canon}}E_C|$. Fact recall and precision are
\begin{equation}
R_{\mathrm{fact}}=\frac{M}{\max(|E_R|,1)},\qquad
P_{\mathrm{fact}}=\frac{M}{\max(|E_C|,1)}.
\end{equation}
SAIS uses the balanced $F_\beta$ form
\begin{equation}
S_{\mathrm{fact}}=
\begin{cases}
1, & |E_R|=|E_C|=0,\\
\frac{(1+\beta_f^2)P_{\mathrm{fact}}R_{\mathrm{fact}}}
{\beta_f^2P_{\mathrm{fact}}+R_{\mathrm{fact}}}, & \text{otherwise},
\end{cases}
\end{equation}
with the score defined as zero if its second-case denominator is zero.
Recall penalizes omitted facts, while precision penalizes extracted
candidate-only facts. Automatic named entity recognition can supply a
baseline set, but curated required facts are recommended for high-stakes use.

\subsection{Constraint Gate}

Facts are not a substitute for requirements. Each constraint $k\in\mathcal K$
has a validation predicate and a failure factor $\gamma_k\in[0,1]$. The gate is
\begin{equation}
\Gamma(C;\mathcal K)=\prod_{k:\operatorname{violated}(C,k)}\gamma_k.
\end{equation}
A hard safety or disclosure violation uses $\gamma_k=0$; a soft discrepancy
can use a factor such as $0.5$. The implementation supports required and
forbidden literal or regular-expression constraints. Domain systems should add
typed validators for structured requirements.

\subsection{Symmetric Length Penalty}

Let $|T|$ be the word count of text $T$. The penalty is
\begin{equation}
BP(C,R)=
\begin{cases}
1, & |C|=|R|=0,\\
0, & |C|=0<|R|,\\
\exp\left(-\beta_\ell\left|\log
\frac{|C|}{\max(|R|,1)}\right|\right), & \text{otherwise},
\end{cases}
\end{equation}
where $\beta_\ell\geq0$. Unlike a one-sided brevity penalty, this form treats
equal proportional shortening and expansion symmetrically. It is a regularizer,
not a proxy for quality, and should be calibrated by task.

\section{Implementation Design}

The accompanying Python package, \texttt{sais-score}, implements the metric
with the following pipeline:
\begin{enumerate}
  \item Accept candidate, reference, optional curated facts, alias map, and
  explicit constraints.
  \item Compute document and block embeddings using a lazy
  \texttt{sentence-transformers} adapter.
  \item Segment blocks, calculate Equation~\ref{eq:logic}, and compute
  document-level Meaning Mesh.
  \item Extract facts using regular expressions and optional spaCy NER, or use
  caller-provided fact lists; canonicalize and match them one-to-one.
  \item Evaluate constraints, apply the length penalty, and return all
  components plus diagnostics.
\end{enumerate}
The package also provides a deterministic hashing encoder for unit tests and
offline demonstrations. It is explicitly not a substitute for a semantic model.

\section{Calibration and Evaluation Protocol}

The hyperparameters $(\alpha,\lambda,\beta_\ell,\beta_f)$ and constraint
factors are task assumptions, not universal constants. They should be selected
on a development partition with no overlap with the final evaluation data.

For empirical validation, use a released human-scored dataset such as a
reference-based WMT Metrics Shared Task set~\citep{freitag2022wmt}. For each
segment, compute SAIS, BLEU, ROUGE-L, and BERTScore under documented
preprocessing. Report Pearson $r$ and Spearman $\rho$ against human scores,
with confidence intervals or bootstrap resampling where possible. Report
language pair, dataset version, segment count, embedding model, entity
extractor, curated-fact policy, and all hyperparameters. The included benchmark
harness accepts local JSONL data and writes a machine-readable result report.

\section{Limitations and Responsible Use}

SAIS inherits embedding bias and can be affected by weak segmentation,
incomplete references, entity extraction errors, and poor alias tables.
Candidate-only named entities are not automatically hallucinations; they may be
valid elaboration. Consequently, the implementation should distinguish curated
critical facts from exploratory automatic extraction. A score must not be used
as the sole basis for medical, legal, employment, educational, safety, or other
high-impact judgments. Human review and task-specific error analysis remain
necessary.

\section{Conclusion}

SAIS separates semantic similarity, logical order, factual precision and
recall, explicit constraints, and length regularization in one interpretable
metric. Its key design choice is to make fact and constraint scores
multiplicative, which prevents high semantic similarity from concealing certain
critical failures. The framework is falsifiable: its value depends on measured
agreement with human judgments and transparent reporting of its calibration.

\bibliographystyle{plainnat}
\bibliography{references}
\end{document}
